% Generated draft: requires mathematical and source review.
\documentclass[11pt]{article}
\usepackage[margin=25mm]{geometry}
\usepackage{fontspec}
\setmainfont{Latin Modern Roman}
\usepackage{amsmath,amssymb,mathtools,unicode-math}
\setmathfont{Latin Modern Math}
\usepackage{enumitem,needspace,xurl,hyperref}
\setlength{\parindent}{0pt}
\setlength{\parskip}{5pt}
\setlength{\emergencystretch}{3em}
\newcommand{\N}{\mathbb{N}}
\newcommand{\Z}{\mathbb{Z}}
\newcommand{\Q}{\mathbb{Q}}
\newcommand{\R}{\mathbb{R}}
\newcommand{\C}{\mathbb{C}}
\newcommand{\F}{\mathbb{F}}
\newcommand{\A}{\mathbb{A}}
\newcommand{\PP}{\mathbb{P}}
\newcommand{\E}{\mathbb{E}}
\newcommand{\eps}{\varepsilon}
\newcommand{\dd}{\,\mathrm{d}}
\begin{document}
\section*{UnsolvedMath ID: 20003129}
\textbf{Source problem code:} catalogue entry
\subsection{Definitions and mathematical statement}
The source asks whether the eventual homeomorphism conclusion of Perelman's stability theorem upgrades to PL homeomorphism or diffeomorphism. A standard precise setting is a sequence of closed connected smooth Riemannian $n$-manifolds $M_i$ with a uniform sectional-curvature lower bound, Gromov--Hausdorff convergence to a compact $n$-dimensional Alexandrov space $X$, and noncollapse (for example a uniform positive volume lower bound in this compact convergence setting). Nonnegative sectional curvature is the workshop's special context. Noncollapse alone, without convergence and the stability hypotheses, does not imply eventual homeomorphism. In the stated stability regime, for some $i_0$ the $M_i$, $i\geq i_0$, are homeomorphic to $X$ and hence to a fixed tail manifold $M$. The upgrade questions are
\[
\exists N\ \forall i,j\geq N,\quad M_i\cong_{\mathrm{PL}}M_j,
\qquad\text{and respectively}\qquad
\exists N\ \forall i,j\geq N,\quad M_i\cong_{\mathrm{diff}}M_j.
\]
The original excerpt omits these surrounding convergence parameters.

The aid gives sufficient topological criteria in dimensions $n\geq5$. Fix homeomorphisms from $M$ to the tail members and pull back their PL structures. Since $M$ already has a smooth structure, a reference PL structure exists. The marked concordance classes of PL structures are controlled by
\[
[M,\operatorname{TOP}/\operatorname{PL}]\cong H^3(M;\Z/2\Z).
\]
Once the PL structure is fixed, smoothing concordance classes are controlled by $[M,\operatorname{PL}/O]$. These are marked/concordance structure targets; they must not be identified directly with unmarked diffeomorphism classes without the relevant concordance and change-of-marking results. Their triviality is a sufficient uniqueness criterion, not a necessary one.

Consequently the supplied criterion is
\[
n\geq5,\quad H^3(M;\Z/2\Z)=0
\quad\Longrightarrow\quad\text{eventual PL homeomorphism},
\]
and, if in addition $[M,\operatorname{PL}/O]$ is the singleton set, eventual diffeomorphism. The aid uses that $\operatorname{PL}/O$ is $6$-connected, giving the explicit special case
\[
n\in\{5,6\},\quad H^3(M;\Z/2\Z)=0
\quad\Longrightarrow\quad\text{eventual diffeomorphism}.
\]
These conclusions depend only on the stabilized topological type and the high-dimensional structure theory; they do not establish general metric control of the structure classes. The same argument does not apply in dimension four.

The aid reports the unrestricted published problem still open, a Pro--Wilhelm result under a codimension-three singular-stratum hypothesis, and an unrefereed 2020 preprint claiming the dimension-four case. No exact statement or full reference is supplied for either special geometric result, so neither is promoted to an unqualified theorem here. In particular the direction and precise meaning of the codimension hypothesis cannot be recovered from the summary alone.
\subsection{Short English statement}
After topological stability, the supplied reduction separates PL structures from smoothings. Vanishing third mod-two cohomology gives a sufficient PL criterion in dimensions at least five, and in dimensions five and six it also gives smooth stability using the connectivity of $\operatorname{PL}/O$. This does not settle the unrestricted problem or dimension four; the reported special geometric results need their full hypotheses and references.
\subsection{Sources}
\begin{itemize}
\item Local registry: permanent catalogue entry, exact source code catalogue entry; \texttt{aim-topology-notes.json}, \texttt{notes[216]}.
\item AIM workshop \emph{Manifolds with nonnegative sectional curvature}, Problem 4, \url{https://aimath.org/WWN/nnsectcurvature/nnsectcurvature.pdf}. Perelman's stability theorem is named, without a full citation, proposer, or proposal year. Original remarks and literature field are empty.
\item Supplied research attempt 1, registry updated 14 August 2026: partial result, low novelty confidence, machine-generated and not peer reviewed. It invokes $\operatorname{TOP}/\operatorname{PL}$ and $\operatorname{PL}/O$ structure theory without complete authoritative references.
\item Its literature report names Pro--Wilhelm's codimension-three singular-stratum case and an unrefereed 2020 preprint claiming dimension four, without titles, theorem numbers, author details for the preprint, or complete citations.
\item The verification field contains only ``The AIM source states:'' and supplies no full stability setup.
\end{itemize}
\end{document}
